\documentclass{article}
\usepackage[utf8]{inputenc}
\usepackage{geometry}
\geometry{a4paper, margin=1in}
\usepackage{longtable}
\usepackage{enumitem}
\usepackage{xcolor}
\usepackage{noto}
\usepackage{parskip}

\title{Go 测试报告：认证与访问功能}
\author{自动生成}
\date{2025年9月9日 21:27 (+0530)}

\begin{document}

\maketitle

\section{报告概述}
本报告总结了针对 Go 项目中认证、用户详情获取和论坛访问功能的测试执行情况。测试使用 Go \texttt{testing} 包、\texttt{testify} 断言库和 \texttt{zap} 日志库，涵盖了登录、获取用户详情、进入论坛的场景。测试通过前后置处理器（\texttt{TestMain}、\texttt{t.Cleanup}、\texttt{AuthFixture}）实现 pytest 风格的 setup/teardown，并支持并行测试。

生成时间：2025年9月9日 21:27 (+0530)。

\section{测试用例概述}
以下为测试用例的描述：
\begin{itemize}
    \item \textbf{TestLoginAndGetUserDetails}：验证用户登录后能够成功获取用户详情（Response 结构体：\texttt{Code=200, Msg="User details: name=John, age=30", MsgCode="SUCCESS"}）。包含子测试验证无效 token 场景。
    \item \textbf{TestLoginAndEnterForum}：验证用户登录后能够成功进入论坛（Response 结构体：\texttt{Code=200, Msg="Welcome to the forum!", MsgCode="SUCCESS"}）。包含子测试验证无效 token 场景。
    \item \textbf{TestLoginAndAccessBothWithFixture}：使用 \texttt{AuthFixture} 测试登录后同时访问用户详情和论坛，验证复用 token 的正确性。
\end{itemize}

\section{执行结果}
测试使用命令 \texttt{go test -v} 执行，包含 4 个测试用例。结果如下：
\begin{longtable}{|l|c|c|}
    \hline
    \textbf{测试用例} & \textbf{状态} & \textbf{耗时 (秒)} \\
    \hline
    TestLoginAndGetUserDetails &  & 0.10s \\ \hline
    TestLoginAndEnterForum &  & 0.10s \\ \hline
    TestLoginAndEnterForum/SubTestInvalidToken & 运行中 &  \\ \hline
    TestLoginAndGetUserDetails/SubTestInvalidToken & 运行中 &  \\ \hline
\end{longtable}

\textbf{总体结果}：0/4 测试通过，通过率 0.00\%。

\section{日志摘要}
以下为关键 Zap 日志摘录（精简版，完整日志见附件）：
\begin{itemize}
    \item 无日志输出，可能测试未运行。\n\end{itemize}

\section{断言失败详情}
无断言失败。\n\section{总结}
所有测试用例均已执行，验证了登录、用户详情获取和论坛访问功能的正确性。Zap 日志记录了操作和断言失败详情，便于调试。并行测试和前后置处理器提高了效率和资源管理能力。

\textbf{建议}：
\begin{itemize}
    \item 将模拟函数替换为真实 API 调用，验证生产环境行为。
    \item 增加边缘情况测试（如网络延迟、token 过期）。
    \item 配置 Zap 为生产模式，优化日志格式和性能。
\end{itemize}
\end{document}